\section{3D Scene Reconstruction}
Building upon previous chapters covering image processing, feature extraction, and multi-view geometry, we now address a new fundamental problem in computer vision:
reconstructing the 3D Euclidean structure of a scene from multiple 2D projections.

\subsection{Structure from Motion (SfM)}
Structure from Motion (SfM) is the computational process of simultaneously estimating the 3D geometry of a scene and the 
camera poses (motion) from a set of 2D images. 

The classical SfM pipeline relies heavily on robust feature correspondence. 
As the camera traverses the scene, perspective transformations and illumination variations alter keypoint appearance, making it hard to establish reliable matches across frames.
To establish reliable correspondences across frames, we extract scale- and rotation-invariant descriptors (such as SIFT, SURF, or ORB) and match them.

\paragraph{Epipolar Geometry and the Essential Matrix}
For any pair of views, the geometric relationship between corresponding points is governed by the epipolar constraint.
Assuming calibrated cameras, let $\mathbf{x}_1$ and $\mathbf{x}_2$ be corresponding points in normalized image coordinates. 
The coplanarity of the camera centers and the 3D point dictates that:
\[
    \mathbf{x}_2^T E \mathbf{x}_1 = 0
\]
where $E$ is the $3 \times 3$ Essential Matrix. 
Algebraically, $E$ encapsulates the relative rotation $R$ and translation $\mathbf{t}$ between the two camera views, defined as:
\[
    E = [\mathbf{t}]_{\times} R
\]

\paragraph{Degrees of Freedom and Estimation}
While $E$ is a $3 \times 3$ matrix containing 9 elements, it possesses strict structural constraints. 
A valid Essential Matrix has exactly 5 degrees of freedom: 
\begin{itemize}
    \item 3 derived from the rotation matrix $R$ 
    \item 2 from the translation vector $\mathbf{t}$ (since the absolute metric scale cannot be recovered from images alone, introducing a scale ambiguity)
\end{itemize}
To estimate $E$, the classical 8-point algorithm provides a linear solution using Singular Value Decomposition (SVD), requiring a minimum of 8 corresponding point pairs. 

\paragraph{Multiple Views}
The two-view approach we have seen so far allows only an initial reconstruction of the scene structure, 
but it is often insufficient for complex scenes.

To achieve a more complete and accurate reconstruction, we can extend the SfM framework to handle multiple views, 
allowing to resolve depth ambiguity and improve reconstruction accuracy.

At this point we have multiple images, each with multiple matched keypoints, but only at a camera pair level.
To fuse all this information together, we can use a global optimization technique called \textbf{bundle adjustment}, 
which jointly refines the 3D point positions and camera parameters to minimize the reprojection error across all views.
This process can be formalized as:
\[
    \arg\min_{\{\mathbf{R}_j, \mathbf{t}_j\}, \{\mathbf{X}_i\}}
    \sum_{i=1}^{N} \sum_{j=1}^{M} 
    v_{ij} \| \mathbf{x}_{ij} - \pi(K, \mathbf{R}_j, \mathbf{t}_j, \mathbf{X}_i) \|^2
\]
where:
\begin{itemize}
    \item $\mathbf{R}_j$ and $\mathbf{t}_j$ are the rotation and translation of the $j$-th camera,
    \item $\mathbf{X}_i$ is the 3D position of the $i$-th keypoint,
    \item $\mathbf{x}_{ij}$ is the observed 2D projection of $\mathbf{X}_i$ in the $j$-th image,
    \item $K$ is the camera intrinsic matrix, and
    \item $v_{ij}$ is a visibility indicator (1 if the point is visible in the image, 0 otherwise).
\end{itemize}

\subsection{Simultaneous Localization and Mapping (SLAM)}
Simultaneous Localization and Mapping (SLAM) is a technique used in robotics and computer vision to build a map of an unknown environment while simultaneously keeping track of the agent's location within that environment.
In SfM, the optimization is performed offline across all views. In the presence of a dense scene and a large number of views, this can be computationally expensive and not suitable for real-time applications.

A usual example of SLAM is the case of a robot navigating in an unknown environment. It needs to estimate its own position while building a map of the environment.
This is a chicken-and-egg problem: to build a map, the robot needs to know where it is; but to know where it is, it needs a map.

\paragraph{The SLAM Pipeline}
Unlike SfM, SLAM systems must operate strictly online and in real-time. The fundamental sequence follows:
\[ \text{Features} \rightarrow \text{Matching} \rightarrow \text{Pose} \rightarrow \text{Map} \]
To achieve this, the pipeline is generally modularized into four core components:
\begin{itemize}
    \item \textbf{Tracking:} Estimates the camera pose continuously as the agent moves.
    \item \textbf{Local Mapping:} Constructs and updates the map of the immediate surroundings.
    \item \textbf{Loop Closure:} Recognizes previously visited areas to correct drift.
    \item \textbf{Global Optimization:} Refines the overall trajectory and map to ensure global consistency.
\end{itemize}

\paragraph{Loop Closure}
In SLAM, the camera is continuously in motion (e.g., mounted on a robot). This means the state of the system (which feeds back into the tracking module) can be defined as:
\[ x_t = f(x_{t-1}, u_t) + w_t \]
Because the system relies on sequential estimations, small errors compound over time, making it necessary to correct the accumulated drift. Recognizing previously visited places helps resolve this issue. Once a loop closure is detected, global optimization can be conducted to minimize the overall re-projection error.

\subsection{SfM vs SLAM}
While they share theoretical foundations in geometry and feature matching, SfM and SLAM are optimized for different operational constraints. 
In particular, SLAM must handle real-time processing and dynamic environments, allowing for online operation and autonomous navigation.
On the other hand, SfM can afford to be more computationally intensive and is typically used for offline reconstruction, resulting in higher accuracy and completeness of the 3D model. 
